\section{Introduction}

The long-only minimum variance portfolio allocates capital across $n$ assets to minimize
portfolio return variance subject to two constraints: the weights must sum to one (the
portfolio is fully invested) and must be non-negative. The only input required in theory
is the $n \times n$ covariance matrix $\Sigma$ of the $n \times 1$ vector of asset
returns. In practice, $\Sigma$ is unknown and must be estimated from historical data.
To this end, one uses a $T\times n$ matrix of demeaned returns $X$, where $T$ denotes
the length of the estimation window. The classical estimator of $\Sigma$ is the sample
covariance matrix $\hat \Sigma = X^\top X/T$.

The Lagrangian for the problem
\[
\min_w\; w^\top \hat \Sigma w \quad \text{subject to} \quad \mathbf{1}^\top w = 1,
\]
where $\mathbf{1}$ denotes a conformable vector of ones, is
\[
\mathcal{L}(w,\lambda) = w^\top \hat \Sigma w - \lambda(\mathbf{1}^\top w - 1),
\]
whose stationarity condition $2 \hat \Sigma w = \lambda\mathbf{1}$ immediately gives
$w \propto \hat \Sigma^{-1}\mathbf{1}$; this was made explicit by \cite{merton1972} in
the context of the efficient frontier. The budget constraint is recovered by the single
rescaling
\[
w = v/(\mathbf{1}^\top v), \qquad v = \hat \Sigma^{-1}\mathbf{1}.
\]
What has \emph{not} been exploited is this structure computationally: the standard
pipeline still forms $\hat \Sigma = X^\top X/T$ explicitly before solving.

This paper makes two contributions. First, we apply conjugate gradients (CG) to
$\hat \Sigma v = \mathbf{1}$ \emph{matrix-free}: $\hat \Sigma = X^\top X/T$ is never
assembled. Each CG iteration requires only two matrix-vector products with $X$, reducing
per-iteration cost from $O(n^2)$ to $O(nT)$ and eliminating $O(n^2)$ storage (the active-set loop further reduces this to $O(kT)$, $k \leq n$; see Table~\ref{tab:complexity}).
Non-negative weights are enforced by a primal-dual active-set loop that alternately removes
assets with negative weight and re-adds assets violating dual feasibility, calling the CG
solver on the modified active set at each step. Because CG is a Krylov method, it builds
successive iterates in the Krylov subspace $\mathcal{K}_k(\hat\Sigma, \mathbf{1})$ and
achieves $O(\sqrt\kappa)$ convergence, where $\kappa = \kappa(\hat\Sigma)$
is the spectral condition number (Proposition~\ref{prop:cg}). This $\sqrt\kappa$ rate is the defining advantage
of Krylov methods over general first-order schemes, which require $O(\kappa)$
iterations without exploiting the subspace structure.

Second, we show that Ledoit-Wolf shrinkage~\cite{ledoit2004}, applied at the heavier
intensities practitioners already use for out-of-sample stability, also compresses the
eigenvalue spread of the system matrix and accelerates CG convergence
(Proposition~\ref{prop:shrinkage}) as a free by-product, with no additional
preconditioning step. On S\&P~500 equity data, heavy LW shrinkage ($\alpha = 0.5$)
reduces CG iterations from 684 to 82; the Frobenius-optimal intensity
$\alpha_{\mathrm{LW}} \approx 0.017$ has a more modest effect (516 iterations), because
the Frobenius loss is insensitive to the near-zero eigenvalues that drive
ill-conditioning. The Frobenius surrogate and the conditioning number are therefore
minimised at different $\alpha$, which might suggest a genuine statistics-versus-speed
trade-off. A rolling-window out-of-sample study (Section~\ref{ssec:oos}) shows there is
none: realised out-of-sample variance is minimised at $\alpha \approx 0.3$, inside the
heavy-shrinkage range that also gives the conditioning benefit, while turnover falls
monotonically in $\alpha$. The intensity that makes the solver fast is thus also the
intensity a risk-conscious practitioner would choose, and the apparent tension is an
artefact of the Frobenius surrogate rather than a property of the portfolio problem.
The observation that heavy shrinkage conditions the solver as a by-product does not
appear to have been made explicit in the portfolio optimisation literature.
RMT eigenvalue cleaning~\cite{laloux1999,bun2017} decouples the two objectives more
decisively; the Woodbury direct solver for the RMT target is developed in the companion
paper~\cite{schmelzer2026rmt}.

CG of Hestenes and Stiefel~\cite{hestenes1952} is a cornerstone of large-scale
scientific computing; see, e.g., Golub and Van Loan~\cite{golub2013}. To the best of
our knowledge it has not been applied to portfolio optimisation in the matrix-free form
described here. The closest related work applies GMRES to the minimum variance problem
when $\hat \Sigma$ is singular~\cite{bajeux2012}; that formulation still assembles
$X^\top X$ and does not exploit the SPD structure or consider preconditioning.
General-purpose solvers such as Clarabel~\cite{clarabel2024} and OSQP~\cite{osqp} handle
the long-only constraint natively but require an assembled covariance matrix as input.
Markowitz's Critical Line Algorithm~\cite{markowitz1956} solves the full
parametric efficient frontier problem via simplex-like pivots ($O(n^2)$ per pivot,
$O(n^3)$ total in the worst case); a clean open-source implementation is available
in~\cite{cvxcla}. The CLA targets the entire analytical frontier rather than a
fixed grid of solves, and does not exploit the matrix-free structure of
$\hat\Sigma = X^\top X/T$.

Computing a discrete efficient frontier via warm-started iterative solves has not,
to the authors' knowledge, been studied as a method in its own right.
The CLA~\cite{markowitz1956,cvxcla} traces the exact analytical frontier but
scales as $O(n^3)$ in the worst case; a na\"{i}ve grid of independent solver calls
does not exploit continuity in the return target $\rho$.
The warm-start described in Section~\ref{sec:results} transfers the active-set mask
and weight vector between adjacent $\rho$ points, so each CG restart begins near
the optimum and the total sweep costs a small multiple of a single solve.
Parametric warm-starting of active-set QP solvers is analysed in the numerical
optimisation literature~\cite{nocedal2006}; combining it with a matrix-free Krylov
inner solver on a changing active set is new.

Three features of the problem make the matrix-free approach non-obvious in standard
QP form. First, all prior solvers treat the assembled $\hat\Sigma$ as a required input;
recognising that $v = \hat\Sigma^{-1}\mathbf{1}$ is the only quantity needed requires
stepping outside the QP abstraction. Second, recovering the budget multiplier
analytically -- reducing the long-only problem to a single SPD linear solve rather than
a constrained QP -- exploits the rank-one structure of the equality constraint, which is
not surfaced by general-purpose QP interfaces. Third, restarting a Krylov inner solver
on a changing active-asset subset each outer step, without re-initialising unnecessarily,
is a non-trivial integration of the active-set and iterative-solver layers.

The scope is deliberately the pure long-only minimum-variance and mean-variance
problem: the matrix-free reduction exploits its specific structure -- a single SPD
solve with a rank-one budget constraint -- and richer constraint sets such as sector
caps, turnover limits, or factor-exposure bounds fall outside it and are better served
by a general-purpose QP solver (Section~\ref{sec:solvers}).

For completeness, Section~\ref{sec:primaldual} also includes a simplex-projection
first-order comparator that handles both constraints simultaneously without an outer active-set loop.

The remainder of the paper is organized as follows.
Sections~\ref{sec:problem}--\ref{sec:pd} develop the problem formulation and
non-negativity strategies. Section~\ref{sec:precond} covers shrinkage targets and preconditioning:
the scaled-identity LW estimator and factor-model targets (companion paper~\cite{schmelzer2026rmt} covers RMT eigenvalue cleaning).
Section~\ref{sec:solvers} presents the concrete solver implementations and
their asymptotic complexity. Section~\ref{sec:results} reports empirical results
on S\&P~500 and FTSE~100 data, a scaling study, and an efficient frontier sweep,
benchmarking against Clarabel, OSQP, and a first-order comparator.
